\documentclass[11pt,a4paper]{article}
\usepackage[spanish,es-noquoting]{babel}
\usepackage[top=2cm, bottom=2cm, left=2.2cm, right=2.2cm]{geometry}
\usepackage{fontspec}
\usepackage{amsmath}
\usepackage[version=4]{mhchem}
\usepackage{enumitem}
\usepackage{microtype}

\setmainfont{Calibri}
\linespread{1.15}

\begin{document}
\pagestyle{empty}

\begin{center}
    {\LARGE \textbf{Resolución: Ejercicios de Rendimiento de Reacciones}}\\[4pt]
    {\large \textbf{Física y Química 4.º ESO / 1.º Bachillerato}}
\end{center}

\vspace{0.3cm}
\noindent\textbf{Concepto clave de rendimiento porcentual:}
\[
\text{Rendimiento (\%)} = \frac{\text{Cantidad real obtenida (experimental)}}{\text{Cantidad teórica calculada por estequiometría}} \times 100
\]

\vspace{0.3cm}
\noindent\textbf{1. Resolución del ejercicio de descomposición de \ce{KClO3}:}
\[
\ce{2 KClO3 (s) -> 2 KCl (s) + 3 O2 (g)}
\]

\noindent\textbf{Cálculo de masas molares ($M$):}
\begin{itemize}[leftmargin=1.5em, itemsep=2pt]
    \item $M(\ce{KClO3}) = 39 + 35{,}5 + 3(16) = 122{,}5\text{ g/mol}$
    \item $M(\ce{O2}) = 2(16) = 32\text{ g/mol}$
\end{itemize}

\begin{enumerate}[label=\textbf{\alph*)}, leftmargin=2em, itemsep=5pt]
    \item \textbf{Moles y masa teórica máxima de \ce{O2} ($100\%$ de rendimiento):}
    \begin{itemize}[leftmargin=1.5em, itemsep=3pt]
        \item Moles de reactivo que reaccionan:
        \[
        n(\ce{KClO3}) = \frac{122{,}5\text{ g}}{122{,}5\text{ g/mol}} = 1{,}00\text{ mol de \ce{KClO3}}
        \]
        \item Según la estequiometría ($2\text{ mol \ce{KClO3}} : 3\text{ mol \ce{O2}}$):
        \[
        n_{\text{teórico}}(\ce{O2}) = 1{,}00\text{ mol \ce{KClO3}} \times \frac{3\text{ mol \ce{O2}}}{2\text{ mol \ce{KClO3}}} = \mathbf{1{,}50\text{ mol de \ce{O2}}}
        \]
        \item Masa teórica máxima producida:
        \[
        m_{\text{teórica}}(\ce{O2}) = 1{,}50\text{ mol} \times 32\text{ g/mol} = \mathbf{48{,}0\text{ g de \ce{O2}}}
        \]
    \end{itemize}

    \item \textbf{Porcentaje de rendimiento del proceso:}
    \[
    \text{Rendimiento} = \frac{m_{\text{real}}}{m_{\text{teórica}}} \times 100 = \frac{38{,}4\text{ g}}{48{,}0\text{ g}} \times 100 = \mathbf{80\%}
    \]
\end{enumerate}

\vspace{0.4cm}
\noindent\textbf{2. Resolución del ejercicio de reacción de \ce{Zn} con \ce{HCl} (Rendimiento = 75\%):}
\[
\ce{Zn (s) + 2 HCl (aq) -> ZnCl2 (aq) + H2 (g)}
\]

\noindent\textbf{Cálculo de masas molares ($M$):}
\begin{itemize}[leftmargin=1.5em, itemsep=2pt]
    \item $M(\ce{Zn}) = 65{,}4\text{ g/mol} \qquad M(\ce{ZnCl2}) = 65{,}4 + 2(35{,}5) = 136{,}4\text{ g/mol}$
    \item Moles de zinc iniciales: $n(\ce{Zn}) = \frac{65{,}4\text{ g}}{65{,}4\text{ g/mol}} = 1{,}00\text{ mol de \ce{Zn}}$.
\end{itemize}

\begin{enumerate}[label=\textbf{\alph*)}, leftmargin=2em, itemsep=5pt]
    \item \textbf{Masa teórica de cloruro de zinc (\ce{ZnCl2}) al $100\%$:}
    \begin{itemize}[leftmargin=1.5em, itemsep=3pt]
        \item Relación estequiométrica $1:1 \implies n_{\text{teórico}}(\ce{ZnCl2}) = 1{,}00\text{ mol de \ce{ZnCl2}}$.
        \item Masa teórica: $m_{\text{teórica}}(\ce{ZnCl2}) = 1{,}00\text{ mol} \times 136{,}4\text{ g/mol} = \mathbf{136{,}4\text{ g de \ce{ZnCl2}}}$.
    \end{itemize}

    \item \textbf{Masa real de \ce{ZnCl2} obtenida (Rendimiento = $75\%$):}
    \[
    m_{\text{real}}(\ce{ZnCl2}) = m_{\text{teórica}} \times 0{,}75 = 136{,}4\text{ g} \times 0{,}75 = \mathbf{102{,}3\text{ g de \ce{ZnCl2}}}
    \]

    \item \textbf{Volumen real de \ce{H2} en Condiciones Normales ($\text{C.N.}$):}
    \begin{itemize}[leftmargin=1.5em, itemsep=3pt]
        \item Moles teóricos de \ce{H2} ($1:1$): $n_{\text{teórico}}(\ce{H2}) = 1{,}00\text{ mol de \ce{H2}}$.
        \item Moles reales producidos al $75\%$: $n_{\text{real}}(\ce{H2}) = 1{,}00\text{ mol} \times 0{,}75 = 0{,}75\text{ mol de \ce{H2}}$.
        \item Volumen real en C.N. ($V_{\text{molar}} = 22{,}4\text{ L/mol}$):
        \[
        V_{\text{real}}(\ce{H2}) = 0{,}75\text{ mol} \times 22{,}4\text{ L/mol} = \mathbf{16{,}8\text{ L de \ce{H2}}}
        \]
    \end{itemize}
\end{enumerate}

\end{document}
